Ce projet a permis de concevoir et comparer plusieurs stratégies d'ordonnancement pour le Serverless Edge Computing : heuristiques (APE-H) du premier semestre et agents d'apprentissage par renforcement (PPO, DQN, Multi-Agent MoE) du second. Les résultats montrent que le \textbf{PPO} offre le meilleur compromis global — il maintient presque toutes les requêtes sous le seuil SLA tout en gardant une consommation énergétique raisonnable — tandis que l'architecture Multi-Agent excelle sur l'économie énergétique au prix d'une latence moins maîtrisée. En perspective, un entraînement conjoint des experts MoE et l'intégration de contraintes dures sur la latence permettraient d'améliorer encore ces résultats.
